\subsection{Minimal refinement}
\label{sec:minimal-refinement}

Once a contradiction shows that the current atom is too coarse, the next question is not how to grow a large model. It is which information must be added. The primitive operation in this paper is minimal refinement: given an information structure and evidence that a new event or variable must be observable, add that observability and nothing else. We state the finite partition theorem first, then the sigma-algebra form. The latter is the standard measure-theoretic join; the former is its finite special case \citep{billingsley1995probability,kallenberg2002foundations}.

\subsubsection{The finite refinement rule}

Let $\calP$ be a finite partition and let $D\subseteq\Omega$ be a new binary distinction event. Define
\begin{align}
\calP_D
&=\{B\cap D:B\in\calP, B\cap D\ne\varnothing\}\notag\\
&{}\cup\{B\cap D^c:B\in\calP, B\cap D^c\ne\varnothing\}.
\label{eq:finite-refinement}
\end{align}
This keeps exactly the nonempty pieces obtained by cutting each old atom by $D$.

\begin{theorem}[Minimal finite refinement]
\label{thm:minimal-refinement:informal}
Let $\calP$ be a finite partition of $\Omega$ and let $D\subseteq\Omega$. The partition $\calP_D$ is the unique minimal partition that refines $\calP$ and makes $D$ measurable. More precisely, if $\calQ$ refines $\calP$ and $D$ is $\calQ$-measurable, then $\calQ$ refines $\calP_D$.
\end{theorem}

\begin{proof}
Each old atom $B\in\calP$ is replaced by the nonempty sets among $B\cap D$ and $B\cap D^c$, so the resulting sets are disjoint, cover $\Omega$, and each lies inside an atom of $\calP$. Hence $\calP_D$ is a refinement of $\calP$. Moreover,
\begin{align*}
D=\bigcup_{B\in\calP}(B\cap D),
\end{align*}
so $D$ is a union of $\calP_D$-atoms. Now let $\calQ$ refine $\calP$ and make $D$ measurable. For every $Q\in\calQ$, there is $B\in\calP$ with $Q\subseteq B$, and measurability of $D$ with respect to $\calQ$ implies either $Q\subseteq D$ or $Q\subseteq D^c$. Thus $Q$ is contained in either $B\cap D$ or $B\cap D^c$, i.e., in an atom of $\calP_D$. Therefore $\calQ$ refines $\calP_D$.
\end{proof}

The theorem gives a canonical update. The evidence determines the least refined representation: atoms untouched by $D$ remain unchanged, while atoms crossed by $D$ split into the two pieces that the evidence distinguishes.

\subsubsection{The sigma-algebra form}

For general spaces, the same operation is the join of sigma-algebras. If the current information is $\calF_t$ and the evidence is a variable $U_t$, the minimal updated information is
\begin{align}
\calF_{t+1}=\calF_t\vee\sigma(U_t).
\label{eq:join-update}
\end{align}

\begin{theorem}[Minimal sigma-algebra refinement]
\label{thm:sigma-refinement:informal}
Let $\calF\subseteq\calE$ be a sigma-algebra and let $U$ be a random variable. The join $\calF\vee\sigma(U)$ is the smallest sigma-algebra $\calG$ such that $\calF\subseteq\calG$ and $U$ is $\calG$-measurable.
\end{theorem}

\begin{proof}
By definition, $\calF\vee\sigma(U)=\sigma(\calF\cup\sigma(U))$, so it contains $\calF$ and makes $U$ measurable. If another sigma-algebra $\calG$ contains $\calF$ and makes $U$ measurable, then $\calF\cup\sigma(U)\subseteq\calG$. Since $\calG$ is a sigma-algebra, $\sigma(\calF\cup\sigma(U))\subseteq\calG$.
\end{proof}

Thus the same object governs finite decision-tree splits and measurable-space refinement. A distinction is not an auxiliary feature added after learning; it is a change in the sigma-algebra of what the learner may condition on. The obstruction in Section~\ref{sec:preliminaries} says when such a change is necessary; minimal refinement says how little information must be added.
